\subsection{Actual coordinate charts of the native polynomial Proj scheme}
\label{subsec:native-projective-charts}

The polynomial ring has its actual total-degree grading over any commutative
coefficient ring. The affine coordinate ring below is the existing
$R[U_i:i\ne j]$. We construct genuine homogeneous-localization and actual
scheme-section comparisons; no chart-algebra certificate is assumed.
This chapter proves the algebraic scheme and section interfaces. It does not
yet assert a global comparison with normalized graded seminorm points.

\begin{definition}[Actual degree-zero chart coordinates]
\label{projective:native-homogeneous-chart-coordinates}
\lean{EnvelopeProperty.ProjBerk.ProjectiveGrading,
  EnvelopeProperty.ProjBerk.HomogeneousChartRing,
  EnvelopeProperty.ProjBerk.homogeneousCoefficientZero,
  EnvelopeProperty.ProjBerk.homogeneousChartCoefficients,
  EnvelopeProperty.ProjBerk.homogeneousChartAlgebra,
  EnvelopeProperty.ProjBerk.homogeneousChartVal,
  EnvelopeProperty.ProjBerk.homogeneousChartCoordinate,
  EnvelopeProperty.ProjBerk.homogeneousChartRatioHom,
  EnvelopeProperty.ProjBerk.homogeneousChartVal_comp_ratioHom}
\leanok
\uses{charts:ratio-map}
For $P=R[X_0,\ldots,X_n]$ with its total-degree grading, let
$P_{(X_j)}$ be Mathlib's genuine degree-zero homogeneous localization
away from $X_j$. Constants give its actual $R$-algebra action. The
map $R[U_i:i\ne j]\to P_{(X_j)}$ sends every $U_i$ to the actual
homogeneous fraction $X_i/X_j$. Its inclusion into the full ordinary
localization is the previously proved whole chart ratio homomorphism.
\end{definition}
\begin{proof}
The actual degree-zero graded piece contains every constant polynomial,
and its coefficient ring homomorphism followed by the homogeneous
localization coefficient inclusion gives the real action. Each coordinate
ratio has numerator and denominator both in degree one, so is an actual
homogeneous-localization element. The polynomial universal property
constructs the map, preserving every polynomial by its ring-map laws.
\end{proof}

\begin{theorem}[Genuine homogeneous chart algebra isomorphism]
\label{projective:native-homogeneous-chart-isomorphism}
\lean{EnvelopeProperty.ProjBerk.coordinateDehomLocalization,
  EnvelopeProperty.ProjBerk.coordinateDehomLocalization_algebraMap,
  EnvelopeProperty.ProjBerk.homogeneousChartDehom,
  EnvelopeProperty.ProjBerk.homogeneousChartDehom_mk,
  EnvelopeProperty.ProjBerk.homogeneousChartVal_injective,
  EnvelopeProperty.ProjBerk.homogeneousChartDehom_comp_ratioHom,
  EnvelopeProperty.ProjBerk.homogeneousChartRatioHom_comp_dehom,
  EnvelopeProperty.ProjBerk.homogeneousChartEquiv}
\leanok
\uses{projective:native-homogeneous-chart-coordinates,charts:homogeneous-substitution}
There is an actual $R$-algebra isomorphism
$P_{(X_j)}\simeq R[U_i:i\ne j]$ sending $F/X_j^d$ to the actual
$X_j=1$ dehomogenization of each degree-$d$ form $F$.
\end{theorem}
\begin{proof}
Ordinary dehomogenization sends $X_j$ to one, so extends through its
ordinary localization. Compose this actual localization homomorphism
with inclusion of the homogeneous degree-zero subring. On the true
coordinate ratios it recovers $U_i$, so the inverse law on the entire
affine polynomial ring follows from its universal property.

Every homogeneous-localization element has an actual representative
$F/X_j^d$ with $F$ homogeneous of degree $d$, by the genuine fraction
surjectivity theorem. The earlier whole homogeneous scaling identity
says that substituting ratios into its dehomogenization, followed by
multiplication by $X_j^d$, gives $F$ in the ordinary localization.
The denominator is a unit, so this proves the inverse identity on every
such fraction. The actual homogeneous-localization inclusion is injective,
yielding the second whole-ring inverse. The argument needs no domain,
field or nonzero coefficient-ring hypothesis.
\end{proof}

\begin{theorem}[Actual affine scheme sections on a coordinate open]
\label{projective:native-scheme-coordinate-sections}
\lean{EnvelopeProperty.ProjBerk.NativeProjectiveScheme,
  EnvelopeProperty.ProjBerk.nativeCoordinateOpen,
  EnvelopeProperty.ProjBerk.nativeHomogeneousChartSchemeIso,
  EnvelopeProperty.ProjBerk.nativeCoordinateOpen_isAffine,
  EnvelopeProperty.ProjBerk.nativeCoordinateSchemeIso,
  EnvelopeProperty.ProjBerk.nativeHomogeneousChartSectionsIso,
  EnvelopeProperty.ProjBerk.nativeCoordinateSectionsIso,
  EnvelopeProperty.ProjBerk.nativeHomogeneousChartSectionsIso_hom,
  EnvelopeProperty.ProjBerk.nativeHomogeneousChartSectionsIso_germ}
\leanok
\uses{projective:native-homogeneous-chart-isomorphism}
The actual scheme $\operatorname{Proj}P$ restricted to the coordinate
open $D_+(X_j)$ is isomorphic to $\operatorname{Spec}R[U_i:i\ne j]$.
Its actual local section ring is genuinely isomorphic to that coordinate
ring, preserving the above homogeneous fractions.
\end{theorem}
\begin{proof}
Mathlib's actual Proj affine-open isomorphism identifies this scheme open
with the spectrum of the genuine homogeneous degree-zero localization,
since $X_j$ has positive degree one. Apply Spec to the actual ring
isomorphism above and compose these genuine scheme isomorphisms.
Their actual global section maps, combined with the Spec global-section
isomorphism and the open-scheme global sections isomorphism, give the
section-ring comparison. This uses the real structure sheaf of Proj.

The defining structure-sheaf compatibility of the actual Proj chart says
that the Spec global-section map followed by the chart pullback equals
Mathlib's homogeneous fraction section homomorphism. The open-scheme
global-section isomorphism cancels its inverse, proving equality of the
section comparison with that fraction homomorphism on the whole ring.
Composing with each genuine germ map then gives the localization map into
the actual projective stalk, using Mathlib's stalk isomorphism. Thus this
comparison preserves every homogeneous fraction and every such germ; it
does not merely compare coordinate values.
\end{proof}
